\documentclass[book.tex]{subfiles}
\begin{document}

\chapter{Differentiable Robotic Simulation}
\label{chap:lagrangian_nn}
\noindent\rule{11cm}{0.4pt} \\
\textbf{Prerequisites:} \autoref{chap:lagrangian_mechanics}, \autoref{chap:hamiltonian_nn} \\
\textbf{Difficulty Level:} ** \\
\noindent\rule{11cm}{0.4pt}
\vspace{5mm}

%Dojo: A Differentiable Simulator for Robotics
Robotics simulation poses an interesting challenge problem for mechanical modeling since it requires high fidelity modeling of friction and contact dynamics. Simulations must also retain fidelity for long time-scales, posing a challenge given the non-differentiable modeling constraints. In this chapter, we study techniques for stabilizing robotic simulations including maximal coordinate representations, smoothing mechanics for contact dynamics, and a cone model for friction dynamics. Implicit differentiable layers provide a unified framework for addressing these questions.

%"""
%We present a differentiable rigid-body-dynamics simulator for robotics that prioritizes physical accuracy and differentiability: Dojo. The simulator utilizes an expressive maximal-coordinates representation, achieves stable simulation at low sample rates, and conserves energy and momentum by employing a variational integrator. Dojo allows for simulation of hard contact that accurately models physical interactions of rigid robotic hardware. In addition, Dojo provides smooth, analytical gradients of contact dynamics, thus bridging the gap between hard and soft contact models by leveraging state-of-the-art interior-point smoothing and continuation strategies. A nonlinear complementarity problem, with nonlinear friction cones, models hard contact and is reliably solved using a custom primal-dual interior-point method. The implicit-function theorem enables efficient differentiation of intermediate results to compute smooth gradients from the contact model. We demonstrate the usefulness of the simulator and its gradients through a number of examples including: simulation, trajectory optimization, reinforcement learning, and system identification."""

Differentiable simulators of physics for robotics systems have become increasingly important for robotic systems \cite{howell2022dojo}.

\section{Differentiable Simulation}

How can we make the simulation operation differentiable for an instantaneous velocity change?

\begin{align}
    v^+ &= v^- + \Delta v
\end{align}

If $\Delta v$ is differentiable, then the update becomes differentiable. More generally, a differentiable contact model can be broken into a series of phases: the maximum dissipation principle, the compression phase, and the restitution phase impulse.

All of these phases can be formulated as convex optimization problems. We can use a differentiable convex layer to implement this.

\begin{align}
\underset{f_C}{\textrm{minimize }} \frac{1}{2}f_C^T A  f_C + f_C v^-
\end{align}

Constrained Lagrangian and Hamiltonian networks can be extended with a differentiable contact model to model a broader class of networks. %Also proposed a regularized version of the contact models.

\begin{align}
(\dot{x}, \dot{v}) &= g(x, v; M(x))
\end{align}
The contact model is given by
\begin{align}
    \Delta v &= \textrm{ContactModel}(x, v; M_{\theta}(x), \mu, e_P) \\
    v&= v' + \Delta v
\end{align}

The framework can learn inertia, potential energy, and contact properties. %Demonstrate on bouncing contact system and bouncing disk system.

%\section{Interpretability}

%\begin{align}
%    L &= \frac{1}{2}m\dot{q}^2 + mg(1-\cos q) \\
%    \ddot{q} &= ..
%\end{align}
%This is the Lagrangian of a simple pendulum. Note mass isn't directly encoded.

%\textcolor{red}{TODO (CITE papers):
%- Incremental Potential Contact
%- DiffTaichi
%- NeuralSim
%- GradSim   
%- ODE Event function
%- Analytically differentiable dynamics}

\section{Maximal Coordinates Representation}
Robotic systems often have closed kinematic loops, which introduce numerical ill-conditioning to simulations.

Minimal coordinate representations use the base $(x,y,z)$ for a base and angular offshifts $(\theta_1, \theta_2)$ for robotic hands. This makes a quadruped robot have 37 dimensions. In a maximal representation, every part of the robot has its own direct representation $(x_i, y_i, z_i)$. This makes the representation dimension much higher.

\section{Variational Integrator}

A variational integrator can be useful for applications where there is good conservation of linear momentum, angular momentum, energy.

\begin{align}
    \mathcal{S} &= \int_{t_1}^{t_2} \mathcal{L} dt
\end{align}
The Euler-Lagrange equations can be applied to transform this in to $F=ma$. The discretization is often applied at the end. A variational integrator instead discretizes the action
\begin{align}
    \mathcal{S}_D &= h \sum_{i=1}^N \mathcal{L}_i
\end{align}

This scheme helps better enforce energy conservation so we can take longer timesteps.

\section{Higher Fidelity Contact Model}

Most contact models have approximations which create unphysical behaviors. For example, approximate friction can cause drift during sliding. Another approximation for impact is soft contact, which can cause impacting objects to sink upon impact and have floor penetration which is aphysical.

It's common to have a "friction cone" in which friction is minimized %(\textcolor{red}{TODO: Cite}). 
Kinetic energy is then minimized within this cone. A common approximation is to use a friction "pyramid" of regions. This approximation is aphysical though and can lead to error terms.

Conic optimization instead can be used to directly handle the nonlinear friction model.

\subsection{Impact Model}

There is a distance constraint $\phi \geq 0$ and the impact force $\gamma \geq 0$. There is a constraint

\begin{align}
    \gamma \circ \phi(x) &= 0
\end{align}

This combination of constraints allows for better simulation of hard contacts. The optimization problem to be solved is given by

\begin{align}
   & \underset{x}{\textrm{minimize}} f(x) \\
   & \textrm{subject to } \phi(x) \geq 0
\end{align}
This can be converted to a soft constraint
\begin{align}
    \textrm{minimize} f(x) - \kappa \log(\phi(x))
\end{align}

The smooth version of the objective allows for better optimization %(\textcolor{red}{TODO: Isn't this just standard?}.) 
As $\kappa \to 0$, the approximation gets steadily better. The KKT conditions are

\begin{align}
   & \textrm{find } x, y \\
   & \textrm{subject to }  \nabla f(x) + \nabla \phi(x)^T \gamma = 0 \\
   & \gamma \cdot \phi(x) = 0 \\
   & \phi(x), \gamma \geq 0
\end{align}

As $\kappa \to 0$, the approximation to hard contact improves. $\kappa$ is iteratively improved through the simulation

\section{Gradient Handling Through Simulator}

Picture a block being lifted off the ground. At start, the block won't lift until the force exceeds gravity, which creates a nonsmooth (ReLu-like) behavior. There is then a discontinuity in the gradient. %People have sample gradients, and there's a new idea of gradient bundles. (\textcolor{red}{TODO: Is this a sample of gradients?}). 
The smooth approximation used in this method allows for smooth gradients.

The actual interior point solver is differentiated through by the implicit function theorem.
\begin{align}
\frac{\partial w^*}{\partial \theta} &= - \left (\frac{\partial r}{\partial w}\right )^{-1} \frac{\partial r}{\partial \theta}
\end{align}
Here $w^*$ is the solution, $\theta$ is the data. %(\textcolor{red}{TODO: We need better material on implicit function theorem.})

%\section{Trajectory Optimization}

%Optimizing trajectories can be performed with iterative LQR. The differentiable simulator can also work for augmented random search in reinforcement learning.

%The simulator could also be used for system identification of the environment.

%\begin{align}
%    \mathcal{L}(D, \theta) &= \sum_{Z \in D} L(Z, \theta) + ...
%\end{align}



\end{document}
